% Options for packages loaded elsewhere
\PassOptionsToPackage{unicode}{hyperref}
\PassOptionsToPackage{hyphens}{url}
\documentclass[
  ignorenonframetext,
]{beamer}
\newif\ifbibliography
\usepackage{pgfpages}
\setbeamertemplate{caption}[numbered]
\setbeamertemplate{caption label separator}{: }
\setbeamercolor{caption name}{fg=normal text.fg}
\beamertemplatenavigationsymbolsempty
% remove section numbering
\setbeamertemplate{part page}{
  \centering
  \begin{beamercolorbox}[sep=16pt,center]{part title}
    \usebeamerfont{part title}\insertpart\par
  \end{beamercolorbox}
}
\setbeamertemplate{section page}{
  \centering
  \begin{beamercolorbox}[sep=12pt,center]{section title}
    \usebeamerfont{section title}\insertsection\par
  \end{beamercolorbox}
}
\setbeamertemplate{subsection page}{
  \centering
  \begin{beamercolorbox}[sep=8pt,center]{subsection title}
    \usebeamerfont{subsection title}\insertsubsection\par
  \end{beamercolorbox}
}
% Prevent slide breaks in the middle of a paragraph
\widowpenalties 1 10000
\raggedbottom
\AtBeginPart{
  \frame{\partpage}
}
\AtBeginSection{
  \ifbibliography
  \else
    \frame{\sectionpage}
  \fi
}
\AtBeginSubsection{
  \frame{\subsectionpage}
}
\usepackage{iftex}
\ifPDFTeX
  \usepackage[T1]{fontenc}
  \usepackage[utf8]{inputenc}
  \usepackage{textcomp} % provide euro and other symbols
\else % if luatex or xetex
  \usepackage{unicode-math} % this also loads fontspec
  \defaultfontfeatures{Scale=MatchLowercase}
  \defaultfontfeatures[\rmfamily]{Ligatures=TeX,Scale=1}
\fi
\usepackage{lmodern}
\ifPDFTeX\else
  % xetex/luatex font selection
\fi
% Use upquote if available, for straight quotes in verbatim environments
\IfFileExists{upquote.sty}{\usepackage{upquote}}{}
\IfFileExists{microtype.sty}{% use microtype if available
  \usepackage[]{microtype}
  \UseMicrotypeSet[protrusion]{basicmath} % disable protrusion for tt fonts
}{}
\makeatletter
\@ifundefined{KOMAClassName}{% if non-KOMA class
  \IfFileExists{parskip.sty}{%
    \usepackage{parskip}
  }{% else
    \setlength{\parindent}{0pt}
    \setlength{\parskip}{6pt plus 2pt minus 1pt}}
}{% if KOMA class
  \KOMAoptions{parskip=half}}
\makeatother
\setlength{\emergencystretch}{3em} % prevent overfull lines
\providecommand{\tightlist}{%
  \setlength{\itemsep}{0pt}\setlength{\parskip}{0pt}}
% Slide layout defaults for warning-clean scientific decks.
\usepackage{etoolbox}
\IfFileExists{xurl.sty}{\usepackage{xurl}}{}
\IfFileExists{seqsplit.sty}{\usepackage{seqsplit}}{\newcommand{\seqsplit}[1]{#1}}
\protected\def\breakseq#1{\seqsplit{#1}}
\protected\def\breaktt#1{\begingroup\ttfamily\seqsplit{#1}\endgroup}
\setlength{\emergencystretch}{6em}
\tolerance=5000
\hbadness=10000
\hfuzz=1pt
\setlength{\tabcolsep}{2pt}
\AtBeginEnvironment{longtable}{\tiny\renewcommand{\arraystretch}{0.86}\setlength{\tabcolsep}{1pt}}
\AtBeginEnvironment{tabular}{\tiny\renewcommand{\arraystretch}{0.86}\setlength{\tabcolsep}{1pt}}

% Natbib and cross-reference fallbacks — slides don't load natbib
% or cleveref, but combined-PDF manuscript prose may emit these
% commands. The fallback renders citations as a bracketed key list
% and cross-references as detokenized labels so slides don't fail on
% undefined control sequences or raw underscores. \providecommand is
% a no-op if packages load later.
\providecommand{\citep}[1]{[#1]}
\providecommand{\citet}[1]{#1}
\providecommand{\citealp}[1]{#1}
\providecommand{\citeauthor}[1]{#1}
\providecommand{\citeyear}[1]{#1}
\providecommand{\cref}[1]{\texttt{\detokenize{#1}}}
\providecommand{\Cref}[1]{\texttt{\detokenize{#1}}}

\newtheorem{proposition}{Proposition}
\newtheorem{hypothesis}{Hypothesis}
\usepackage{bookmark}
\IfFileExists{xurl.sty}{\usepackage{xurl}}{} % add URL line breaks if available
\urlstyle{same}
% fallback for those not using the hyperref driver hyperxmp:
\makeatletter
\@ifundefined{xmpquote}{\newcommand{\xmpquote}[1]{#1}}{}
\makeatother
\hypersetup{
  hidelinks,
  pdfcreator={LaTeX via pandoc}}

\author{\texorpdfstring{}{}}
\date{}

\begin{document}

\section{Case Studies: Dementia Care as a Stress
Test}\label{sec:casestudies}

\begin{frame}[allowframebreaks]{Case Studies: Dementia Care as a Stress
Test}
We illustrate the framework with one worked case study in dementia care.
It is presented not as evidence of efficacy but as a boundary test: a
setting chosen because it strains several of DigiPPPiP's core
assumptions rather than confirming them. Other applications
(early-childhood dyads, long-distance couples, stroke rehabilitation)
belong to the empirical program in \breakseq{[@sec:agenda]}; here we develop
the single case that most sharply exposes what the kernel does and does
not guarantee.
\end{frame}

\begin{frame}[allowframebreaks]{Dementia Care: Asymmetric Dyads,
Biographical Scaffolding, and Fluctuating Consent}
\protect\phantomsection\label{dementia-care-asymmetric-dyads-biographical-scaffolding-and-fluctuating-consent}
DigiPPPiP's minimum viable kernel (two partners, a shared mark field,
perceptible traces of agency, and consentful control over persistence)
maps onto dementia care dyads in ways that productively stress-test
several of the framework's core assumptions. The most important
asymmetry is cognitive: where the framework presupposes two agents with
roughly stable capacity to govern the shared archive, dementia
introduces radical temporal asymmetry. The adult child experiences each
session as one point in a longitudinal caregiving relationship; the
parent may experience each session as the first. This does not
invalidate the practice. There is substantial evidence that creative
engagement can produce genuine moments of co-presence and recognition
even when sustained narrative continuity is lost; however, it requires
the consent and governance model to be substantially extended. Archive
control cannot rest symmetrically with both partners when one partner's
capacity fluctuates within and across sessions. The framework's default
rule (when partners disagree, the less-disclosing option governs) is a
useful starting point, but dementia care contexts require an explicit
supported decision-making layer, drawing on bioethics literature on
fluctuating consent capacity, that specifies who holds proxy governance,
under what conditions, and with what reversibility.

The cyberphysical substrate takes on new meaning in this context. Rather
than treating the shared canvas as a blank field for improvised
mark-making, a dementia-oriented implementation can draw on biographical
materials (family photographs, familiar places, recurring objects,
handwriting samples) as persistent prompts embedded in the canvas. This
is not merely an accessibility feature but a theoretically motivated
reframing of what the shared surface is. Long-term autobiographical
memory is often more intact than short-term memory in dementia, which
means that a canvas seeded with recognizable materials from the parent's
past can function as a bridge between the child's present and the
parent's accessible memory. The place-based micropractice framing is
especially relevant here: the canvas becomes a relational micro-place
constructed from biographical geography. A kitchen from forty years ago,
a recurring holiday destination, or a childhood home can all serve as
orienting triggers. Recurrence and shared uptake remain the minimum
place claim; in this context they are satisfied when a familiar image,
prompt, or outline reliably orients the parent toward a recognizable
relational moment.

The temporal taxonomy requires reinterpretation rather than extension.
Synchronous co-present sessions remain possible and valuable,
particularly in early-stage dementia, but the semisynchronous and
asynchronous modes acquire specific clinical relevance. An asynchronous
canvas that the child adds to between visits, and that the parent
encounters during a session with a care facilitator, preserves
relational continuity across the gaps that dementia and caregiving
schedules impose. The canvas in this mode functions less as a real-time
communication channel and more as a persistent relational object that
accumulates meaning across time. This is closer to a shared memory book
than a messaging platform, but with the agency and mark-making
properties that distinguish DigiPPPiP from passive reminiscence
materials. The temporal primitive of repair acquires particular weight:
moments where the parent reorients toward the shared surface after
confusion, or where the child responds to an unexpected mark, carry
relational significance that the event logging schema should be designed
to capture without pathologizing disorientation as failure.

The neuroergonomic design constraint is sharpened considerably. The
original principle that the interface should be nearly invisible becomes
a clinical requirement in dementia contexts where cognitive load
management is not a design preference but a condition of participation.
Tool palettes, notifications, and navigational complexity are
contraindicated. The minimum viable interface is a single shared surface
with one input mode, familiar visual prompts, and no decisions required
beyond the mark itself. Accessibility here is not only about motor input
flexibility but about cognitive accessibility: plain-language prompts,
low-distraction environments, and session structures short enough to
remain within the participant's attention window. The care facilitator
becomes a legitimate third role in this context, supporting
participation without substituting for it, and the governance model
should specify the facilitator's access to the canvas and event log
separately from both partners.

The relational aesthetics claim that the artifact points back to the
other person as an agent is both the most fragile and the most important
aspect of this application. A canvas full of marks is not relational if
one partner is merely producing output without awareness of the other's
presence. The honest boundary condition is that DigiPPPiP in dementia
care settings cannot guarantee sustained mutual recognition, and should
not be positioned as doing so. What it can offer is a structured
occasion for moments of recognition: brief, genuine, and worth creating.
Within a practice that asks very little in return, that is a modest
claim. In the context of a condition that progressively forecloses
relational possibility, it is not a small one.
\end{frame}

\end{document}
